% GENERATED FILE. Edit canonical lessons, then run npm run generate:books.
% canonical-source-hash: fnv1a64:c728271089eb6d24
% canonical-lessons: FA-C416-karbar, FA-C416-ramz-e-obur, FA-C416-zakhire-kardan, FA-C416-hazf-kardan, FA-C416-be-ruz-resani

\chapter{Accounts and passwords}
\label{ch:fa-accounts-and-passwords}
\hlchaptermodality{416}

\begin{chapteropening}
\textbf{By the end of this chapter:} \emph{I can talk about users and passwords, saving and deleting things, and updates.}

Complete the last lesson of chapter 416: I can talk about users and passwords, saving and deleting things, and updates.
\end{chapteropening}

\section[\texorpdfstring{\fa{کاربر} (kârbar)}{کاربر (kârbar)}]{\fa{کاربر} (kârbar) — a user}

\label{lesson:FA-C416-karbar}



\begin{quote}
Before the new one: say the Persian for a percent, then the Persian for a spendthrift.
\end{quote}

\subsection*{You'll want to know: \fa{کاربر}}
\textbf{\fa{کاربر}} (\emph{kârbar}) — \textquotedblleft{}a user\textquotedblright{}.

\begin{quote}
\fa{این} \fa{برنامه} \fa{میلیونها} \fa{کاربر} \fa{دارد}.
\end{quote}

\begin{quote}
\emph{in barnâme milyunhâ kârbar dârad.}
\end{quote}

\begin{quote}
This app has millions of users.
\end{quote}

\emph{Kâr} \textquotedblleft{}work\textquotedblright{} + \emph{bar} \textquotedblleft{}carrying\textquotedblright{}. \emph{Nâm-e kârbari} — a username.

\begin{cousinweb}
\textbf{\fa{کاربر}} is a modern native coinage built from \emph{kâr} \textquotedblleft{}work\textquotedblright{} + \emph{bar} \textquotedblleft{}carrying\textquotedblright{}, from \emph{bar-}, the present stem of \emph{bordan} \textquotedblleft{}to carry\textquotedblright{}. A user is one who puts a tool to work.
\end{cousinweb}

\begin{grammarlens}[title={What you've built}]
One more word for telling, explaining and describing at B1.
\end{grammarlens}

\subsection*{Your turn}
\begin{itemize}
\raggedright
  \item \emph{Say it:} \emph{kârbar}
  \item \emph{Say it:} \emph{kârbar}, once more
  \item \emph{Say it:} a sentence of your own with \emph{kârbar}
  \item \emph{Recall:} say the Persian for a percent, then the Persian for a spendthrift, then say \emph{kârbar} again
\end{itemize}

\begin{tcolorbox}[breakable,colback=teal!4,colframe=teal!35!black,title={Before you move on}]
What does \fa{کاربر} mean? (\textquotedblleft{}A user\textquotedblright{}.) Say it once more.
\end{tcolorbox}
\section[\texorpdfstring{\fa{رمز} \fa{عبور} (ramz-e 'obur)}{رمز عبور (ramz-e 'obur)}]{\fa{رمز} \fa{عبور} (ramz-e 'obur) — a password}

\label{lesson:FA-C416-ramz-e-obur}



\begin{quote}
Before the new one: say the Persian for a spendthrift, then the Persian for a user.
\end{quote}

\subsection*{You'll want to know: \fa{رمز} \fa{عبور}}
\textbf{\fa{رمز} \fa{عبور}} (\emph{ramz-e 'obur}) — \textquotedblleft{}a password\textquotedblright{}.

\begin{quote}
\fa{رمز} \fa{عبورم} \fa{را} \fa{فراموش} \fa{کردم}.
\end{quote}

\begin{quote}
\emph{ramz-e 'oburam râ farâmush kardam.}
\end{quote}

\begin{quote}
I forgot my password.
\end{quote}

\emph{Ramz} \textquotedblleft{}code\textquotedblright{} + \emph{'obur} \textquotedblleft{}passing\textquotedblright{}. Also \emph{gozarvâzhe}.

\begin{cousinweb}
\textbf{\fa{رمز} \fa{عبور}} is literally \textquotedblleft{}the code for passing\textquotedblright{}: \emph{ramz} was borrowed from Arabic \emph{ramz} \textquotedblleft{}cypher, code\textquotedblright{} (EtymDB), and \emph{'obur} is usually traced to Arabic \emph{'ubūr} \textquotedblleft{}crossing\textquotedblright{}.
\end{cousinweb}

\begin{grammarlens}[title={What you've built}]
One more word for telling, explaining and describing at B1.
\end{grammarlens}

\subsection*{Your turn}
\begin{itemize}
\raggedright
  \item \emph{Say it:} \emph{ramz-e 'obur}
  \item \emph{Say it:} \emph{ramz-e 'obur}, once more
  \item \emph{Say it:} a sentence of your own with \emph{ramz-e 'obur}
  \item \emph{Recall:} say the Persian for a spendthrift, then the Persian for a user, then say \emph{ramz-e 'obur} again
\end{itemize}

\begin{tcolorbox}[breakable,colback=teal!4,colframe=teal!35!black,title={Before you move on}]
What does \fa{رمز} \fa{عبور} mean? (\textquotedblleft{}A password\textquotedblright{}.) Say it once more.
\end{tcolorbox}
\section[\texorpdfstring{\fa{ذخیره} \fa{کردن} (zakhire kardan)}{ذخیره کردن (zakhire kardan)}]{\fa{ذخیره} \fa{کردن} (zakhire kardan) — to save, to store}

\label{lesson:FA-C416-zakhire-kardan}



\begin{quote}
Before the new one: say the Persian for a user, then the Persian for a password.
\end{quote}

\subsection*{You'll want to know: \fa{ذخیره} \fa{کردن}}
\textbf{\fa{ذخیره} \fa{کردن}} (\emph{zakhire kardan}) — \textquotedblleft{}to save, to store\textquotedblright{}.

\begin{quote}
\fa{عکسها} \fa{را} \fa{در} \fa{کامپیوتر} \fa{ذخیره} \fa{کردم}.
\end{quote}

\begin{quote}
\emph{'akshâ râ dar kâmpyuter zakhire kardam.}
\end{quote}

\begin{quote}
I saved the photos on the computer.
\end{quote}

\emph{Zakhire} \textquotedblleft{}store, reserve\textquotedblright{} + \emph{kardan}. Also to store food or water.

\begin{cousinweb}
\textbf{\fa{ذخیره}} was borrowed from Arabic \emph{dhakhīra} \textquotedblleft{}store, reserve\textquotedblright{} (EtymDB); Persian says the Arabic \emph{dh} as \emph{z}.
\end{cousinweb}

\begin{grammarlens}[title={What you've built}]
One more word for telling, explaining and describing at B1.
\end{grammarlens}

\subsection*{Your turn}
\begin{itemize}
\raggedright
  \item \emph{Say it:} \emph{zakhire kardan}
  \item \emph{Say it:} \emph{zakhire kardan}, once more
  \item \emph{Say it:} a sentence of your own with \emph{zakhire kardan}
  \item \emph{Recall:} say the Persian for a user, then the Persian for a password, then say \emph{zakhire kardan} again
\end{itemize}

\begin{tcolorbox}[breakable,colback=teal!4,colframe=teal!35!black,title={Before you move on}]
What does \fa{ذخیره} \fa{کردن} mean? (\textquotedblleft{}To save, to store\textquotedblright{}.) Say it once more.
\end{tcolorbox}
\section[\texorpdfstring{\fa{حذف} \fa{کردن} (hazf kardan)}{حذف کردن (hazf kardan)}]{\fa{حذف} \fa{کردن} (hazf kardan) — to delete, to remove}

\label{lesson:FA-C416-hazf-kardan}



\begin{quote}
Before the new one: say the Persian for a password, then the Persian for to save.
\end{quote}

\subsection*{You'll want to know: \fa{حذف} \fa{کردن}}
\textbf{\fa{حذف} \fa{کردن}} (\emph{hazf kardan}) — \textquotedblleft{}to delete, to remove\textquotedblright{}.

\begin{quote}
\fa{او} \fa{پیام} \fa{را} \fa{حذف} \fa{کرد}.
\end{quote}

\begin{quote}
\emph{u payâm râ hazf kard.}
\end{quote}

\begin{quote}
She deleted the message.
\end{quote}

\emph{Hazf} \textquotedblleft{}removal\textquotedblright{} + \emph{kardan}. Also to drop a course or a player.

\begin{cousinweb}
\textbf{\fa{حذف}} is usually traced to Arabic \emph{ḥadhf} \textquotedblleft{}cutting off, dropping\textquotedblright{}; Persian says the Arabic \emph{dh} as \emph{z}.
\end{cousinweb}

\begin{grammarlens}[title={What you've built}]
One more word for telling, explaining and describing at B1.
\end{grammarlens}

\subsection*{Your turn}
\begin{itemize}
\raggedright
  \item \emph{Say it:} \emph{hazf kardan}
  \item \emph{Say it:} \emph{hazf kardan}, once more
  \item \emph{Say it:} a sentence of your own with \emph{hazf kardan}
  \item \emph{Recall:} say the Persian for a password, then the Persian for to save, then say \emph{hazf kardan} again
\end{itemize}

\begin{tcolorbox}[breakable,colback=teal!4,colframe=teal!35!black,title={Before you move on}]
What does \fa{حذف} \fa{کردن} mean? (\textquotedblleft{}To delete, to remove\textquotedblright{}.) Say it once more.
\end{tcolorbox}
\section[\texorpdfstring{\fa{به} \fa{روز} \fa{رسانی} (be ruz resâni)}{به روز رسانی (be ruz resâni)}]{\fa{به} \fa{روز} \fa{رسانی} (be ruz resâni) — an update}

\label{lesson:FA-C416-be-ruz-resani}



\begin{quote}
Before the new one: say the Persian for to save, then the Persian for to delete.
\end{quote}

\subsection*{You'll want to know: \fa{به} \fa{روز} \fa{رسانی}}
\textbf{\fa{به} \fa{روز} \fa{رسانی}} (\emph{be ruz resâni}) — \textquotedblleft{}an update\textquotedblright{}.

\begin{quote}
\fa{کامپیوتر} \fa{من} \fa{یک} \fa{به} \fa{روز} \fa{رسانی} \fa{تازه} \fa{گرفت}.
\end{quote}

\begin{quote}
\emph{kâmpyuter-e man yek be ruz resâni-ye tâze gereft.}
\end{quote}

\begin{quote}
My computer got a new update.
\end{quote}

\emph{Be ruz} \textquotedblleft{}up to date\textquotedblright{} + \emph{resâni} \textquotedblleft{}bringing\textquotedblright{}, from \emph{resândan}. \emph{Be ruz resâni kardan} — to update.

\begin{cousinweb}
\textbf{\fa{به} \fa{روز} \fa{رسانی}} is a modern native coinage, literally \textquotedblleft{}bringing to the day\textquotedblright{}: \emph{ruz} \textquotedblleft{}day\textquotedblright{} is traced by EtymDB through Old Persian, and \emph{resâni} is built on \emph{resân-}, the causative stem of \emph{residan} \textquotedblleft{}to arrive\textquotedblright{}.
\end{cousinweb}

\begin{grammarlens}[title={What you've built}]
One more word for telling, explaining and describing at B1.
\end{grammarlens}

\subsection*{Your turn}
\begin{itemize}
\raggedright
  \item \emph{Say it:} \emph{be ruz resâni}
  \item \emph{Say it:} \emph{be ruz resâni}, once more
  \item \emph{Say it:} a sentence of your own with \emph{be ruz resâni}
  \item \emph{Recall:} say the Persian for to save, then the Persian for to delete, then say \emph{be ruz resâni} again
\end{itemize}

\begin{tcolorbox}[breakable,colback=teal!4,colframe=teal!35!black,title={Before you move on}]
What does \fa{به} \fa{روز} \fa{رسانی} mean? (\textquotedblleft{}An update\textquotedblright{}.) Say it once more.
\end{tcolorbox}
